\begin{abstract}
Tool-using language-model agents are hard to trust from a final reply alone. A support bot can sound helpful while skipping authentication, calling the wrong tool, or leaving application state incorrect. Existing benchmarks and evaluation libraries grade outputs and metrics well, but they often split dialogue, tool checks, and state assertions across separate files and return coarse failure signals when something breaks.

This project presents \emph{agent-spec-kit}, a framework for scenario-driven evaluation of LLM agents. Developers write scenarios as readable conversation scripts: user turns, checks on replies and tool traces after each step, and optional state postconditions in one place. Failed assertions produce structured counterexamples that name the broken expectation, the turn, and the field that differed. The same model supports repeated runs, simulated users, mutation fuzzing, failure shrinking, and regression extraction so exploratory bugs can re-enter the suite as permanent tests.

The framework is evaluated on TelcoSupportBench, a synthetic mobile customer-support domain with fifty tasks and a reference LangGraph agent. Studies compare \emph{agent-spec-kit} against seven other approaches, including Ragas and DeepEval, on authoring clarity, failure-message quality, fault detection across oracle lanes, and generative discovery. Results show that multi-surface checks surface failures output-only grading misses; scenario checks are clearer to read and more precise on failure than metric-oriented reimplementations; full oracles catch most injected tool and state faults; and simulation plus fuzzing find many distinct failure modes fixed scripts do not reach. The project argues that dependable agent development needs declarative behavioural specs teams can run, debug, and extend, and that \emph{agent-spec-kit} is a promising step in that direction.
\end{abstract}
